Рассмотрим квадратичную форму вида:
$$ K=\sum_{i,j=1}^n a_{ij}(x^{(0)})y_i y_j = y^T A(x^{(0)})y, \qquad y = (y_1, \dots , y_n) $$
при ЛНЗНП:
\[
y=S\eta,\ \eta=(\eta_1, \ldots, \eta_n)^T,\ detS \neq 0:
\quad
K=y^TA(x^{(0)})y= \eta^T S^T A(x^{(0)}) S\eta =\eta^T\tilde A(\xi^{(0)})\eta
\]

Линейной невырожденной замене независимых переменных

\[y=S\eta\quad (y_i=\sum_{k=1}^ns_{ik}\eta_k)\] соответствует: \[\xi=S^T(x-x^{(0)});\qquad (\xi_k=\sum_{k=1}^ns_{ik}(x_i-x_i^{(0)})) \]

Распишем это подробнее:

Для любой квадратичной формы $K=y^TA(x^{(0)})y$ существует ЛНЗНП: $y=S\eta$, приводящая эту форму к каноническому виду:
\[
  K = \sum_{k=1}^p \eta_k^2 - \sum_{k=p+1}^{p+q} \eta_k^2,
\] 
где $p$ -- ПИИ, $q$ -- ОИИ(индексы инерции); $r=p+q$ -- ранг.

\subsection{Классификация ДУЧП2ПЛОСП в точке}
\begin{enumerate}[label=\Roman*.]
  \item \textbf{Эллиптический тип}
    \begin{itemize}[nosep]
      \item $r = n$ --- ранг $K$
      \item в каноническом виде все слагаемые одного знака 
            ($p = r = n,\ q = 0$ или $q = r = n,\ p = 0$)
      \item \textbf{\underline{Пример:}} уравнение Лапласа.
    \end{itemize}

  \item \textbf{Гиперболический тип}
    \begin{itemize}[nosep]
      \item $r = n$ --- ранг $K$
      \item в каноническом виде слагаемые имеют разные знаки 
            ($p \in [1, n-1]$ и $q \in [1, n-1]$)
      \item \textbf{Замечание:} при $p = n-1,\ q = 1$ или $p = 1,\ q = n-1$ --- 
            \emph{нормальный гиперболический} тип.
      \item \textbf{\underline{Пример:}} волновое уравнение.
    \end{itemize}

  \item \textbf{Параболический тип}
    \begin{itemize}[nosep]
      \item $r < n$ --- ранг $K$
      \item \textbf{Замечание:} при $p < n,\ q = 0$ или $p = 0,\ q < n$ --- 
            \emph{нормальный параболический} тип.
      \item \textbf{\underline{Пример:}} уравнение теплопроводности.
    \end{itemize}
\end{enumerate}

\begin{note}
Ранг $r$, ПИИ --- $p$, ОИИ --- $q$ не инвариантны для фиксированной точки $x^{(0)}$ (или функции фиксированной точки), поэтому тип ДУЧП2ПЛОСП зависит от выбора $x^{(0)} \in G$.
\end{note}

\textbf{\underline{Пример:}} уравнение Трикоми

\[
y \frac{\partial^2 u}{\partial x^2} + \frac{\partial^2 u}{\partial y^2} = 0
\]
Оно является \textbf{эллиптическим} при $y > 0$, \textbf{гиперболическим} при $y < 0$ и \textbf{параболическое} на линии вырождения $y = 0$.

\textbf{\underline{Пример:}} стационарное линеаризованное уравнение для потенциала скорости
\[
(1-M^2) \frac{\partial^2 u}{\partial x^2} + \frac{\partial^2 u}{\partial y^2} + \frac{\partial^2 u}{\partial z^2} = 0
\]
\begin{enumerate}[nosep]
  \item если \(M<1\) (дозвуковое течение), уравнение становится \textbf{эллиптическим};
  \item если \(M>1\) (сверхзвуковое течение), уравнение становится \textbf{гиперболическим};
  \item если \(M=1\) (звуковое течение), уравнение вырождается в \textbf{параболическое}.
\end{enumerate}

\begin{note}
Если в ДУЧП2ПЛОСП старшие коэффициенты $a_{ij} \neq a_{ij}(x)$(т.е. постоянные коэффициенты), то ЛНЗНП $\xi=\xi(x)$ можно привести это уравнение к уравнению канонического вида во всей $G$.
\end{note}

\begin{note}
Лапласиан в декартовой, сферической и цилиндрической системах координат.
\begin{enumerate}
% Прямоугольная
\item Декартовы координаты: $\displaystyle\Delta = \frac{\partial^2}{\partial x^2} + \frac{\partial^2}{\partial y^2} + \frac{\partial^2}{\partial z^2} 
$

% Цилиндрическая
\item Цилиндрические координаты: $\displaystyle\Delta = \frac{1}{r}\frac{\partial}{\partial r}\left(r \frac{\partial}{\partial r}\right) + \frac{1}{r^2}\frac{\partial^2}{\partial \varphi^2} + \frac{\partial^2}{\partial z^2}
$

% Сферическая
\item Сферические координаты: $\displaystyle\Delta = \frac{1}{\rho^2}\frac{\partial}{\partial \rho}\left(\rho^2 \frac{\partial}{\partial \rho}\right) + \frac{1}{\rho^2 \sin\theta}\frac{\partial}{\partial \theta}\left(\sin\theta \frac{\partial}{\partial \theta}\right) + \frac{1}{\rho^2 \sin^2\theta}\frac{\partial^2}{\partial \varphi^2}
$
\end{enumerate}
\end{note}
%FIXME НЕТ КАРТИИИИНОК
Попробуем обнулить коэффициенты $\tilde a_{kl}(\xi)$

\subsection{Характеристические поверхности(линии). Характеристики.}
\begin{definition}[Характеристические поверхности(линии). Характеристики.]
Пусть функция $\omega(x_1, \ldots, x_n) \in C^1(G)$ такова, что на плоскости $\omega(x)=0: \operatorname{grad}(\omega)  \neq0$.
Если на данной плоскости $\omega(x)$ удовлетворяет уравнению 
$\sum_{i, j=1}^n a_{ij}(x)\frac{\partial\omega}{\partial x_i}\frac{\partial\omega}{\partial x_j}=0$, то эта плоскость ($\omega(x)=0$) называется характеристической плоскостью (линией) или характеристикой для ДУЧП2ПЛОСП(1).
\end{definition}

Рассмотрим семейство характеристик $\omega(x, y)=C$, где $\alpha < C < \beta \Rightarrow$ некоторая достаточно малая область $G_\varepsilon$ будет целиком заполнена этими характеристиками, которые не будут пересекаться друг с другом, т.е. через каждую точку из $G_\varepsilon$ будет проходить единственная характеристика.

Если в НЗНП положить $\xi_1=\omega(x)$, то в $(\widetilde{1})$ коэффициент $\tilde a_{11} \equiv 0$.

Рассмотрим возможность приведения ДУЧП2ПЛОСП к каноническому виду в некоторой достаточно малой области $G_\varepsilon$. Для этого необходимо выполнение в $G_\varepsilon$ следующих условий:

\begin{quote}
$\tilde a_{kl}(\xi)=0,\quad k \neq l$

$\tilde a_{kk}(\xi)=\varepsilon_k\tilde a_{11}(\xi),\quad k=\overline{2, n},\ \varepsilon_k=0,\pm1$
\end{quote}

Всего наших соотношений:

\[
N=\frac{n^2-n}{2}+n-1 \leq n,
\]
где $n$ --- число независимых переменных.

Т.е. $n^2-n-2 \leq 0 \Rightarrow (n+1)(n-2) \leq0$, т.е. $n\leq2$

\textbf{\underline{Вывод:}} только лишь при $n \leq 2$ ДУЧП2ПЛОСП может быть приведено к каноническому виду в достаточно малой области с помощью НЗНП.

Рассмотрим случай $\boldsymbol{n=2}$.

\subsection{ДУЧП2ПЛОСП с двумя независимыми переменными}

Итак, \[a(x, y)\frac{\partial^2u}{\partial x^2} + 2b(x, y)\frac{\partial^2u}{\partial x \partial y} + c(x, y)\frac{\partial^2u}{\partial y^2} + \Phi(x, y, u, u_x, u_y)=0;\ (x, y) \in G \subset \mathbb{R}^2, a, b, c \in C(G)\ (1)\]

Рассмотрим НЗНП:
\begin{quote}
$\xi = \xi(x, y) \in C^2(G)$

$\eta = \eta(x, y) \in C^2(G)$
\end{quote}

Тогда в некоторой окрестности точки $\exists$ обратная замена \[x=x(\xi, \eta) \text{ и } y=y(\xi, \eta) \Rightarrow \tilde u(\xi, \eta)=u(x(\xi, \eta), y(\xi, \eta)) \text{ или } u(x, y)=u(\xi(x, y), \eta(x, y))\]

\noindent
Очевидно,
\begin{itemize}
  \item $\displaystyle 
    \frac{\partial u}{\partial x}
    = \frac{\partial \tilde u}{\partial \xi}\xi_x
    + \frac{\partial \tilde u}{\partial \eta}\eta_x
  $
  \item $\displaystyle 
    \frac{\partial u}{\partial y}
    = \frac{\partial \tilde u}{\partial \xi}\xi_y
    + \frac{\partial \tilde u}{\partial \eta}\eta_y
  $
  \item $\displaystyle 
    \frac{\partial^2 u}{\partial x^2}
    = \frac{\partial^2 \tilde u}{\partial \xi^2}\xi_x^2
    + 2\frac{\partial^2 \tilde u}{\partial \xi \partial \eta}\xi_x\eta_x
    + \frac{\partial^2 \tilde u}{\partial \eta^2}\eta_x^2
    + \frac{\partial \tilde u}{\partial \xi}\xi_{xx}
    + \frac{\partial \tilde u}{\partial \eta}\eta_{xx}
  $
  \item $\displaystyle 
    \frac{\partial^2 u}{\partial x \partial y}
    = \frac{\partial^2 \tilde u}{\partial \xi^2}\xi_x\xi_y
    + \frac{\partial^2 \tilde u}{\partial \xi \partial \eta}(\xi_x\eta_y + \xi_y\eta_x)
    + \frac{\partial^2 \tilde u}{\partial \eta^2}\eta_x\eta_y
    + \frac{\partial \tilde u}{\partial \xi}\xi_{xy}
    + \frac{\partial \tilde u}{\partial \eta}\eta_{xy}
  $
  \item $\displaystyle 
    \frac{\partial^2 u}{\partial y^2}
    = \frac{\partial^2 \tilde u}{\partial \xi^2}\xi_y^2
    + 2\frac{\partial^2 \tilde u}{\partial \xi \partial \eta}\xi_y\eta_y
    + \frac{\partial^2 \tilde u}{\partial \eta^2}\eta_y^2
    + \frac{\partial \tilde u}{\partial \xi}\xi_{yy}
    + \frac{\partial \tilde u}{\partial \eta}\eta_{yy}
  $
\end{itemize}
